\section{Panel-aware generative representation learning}

We developed a panel-aware generative model to learn cell embeddings from single-cell
expression profiles measured over heterogeneous gene panels. Each cell is represented
by a count vector \(x_i \in \mathbb{N}^{G}\) and a binary observation mask
\(m_i \in \{0,1\}^{G}\), where \(m_{ig}=1\) indicates that gene \(g\) was measured in
cell \(i\). Counts for unobserved genes are not treated as zeros in the likelihood.
Instead, the mask is explicitly provided to the encoder and reconstruction losses are
computed only on observed entries unless otherwise specified.

\subsection{Mask-conditioned encoder}

The encoder receives log-transformed counts and the corresponding observation mask.
Specifically, we encode \(\log(1+x_i)\) with an expression branch and \(m_i\) with a
mask branch. The two branches are fused by a feature-wise affine modulation:
\begin{align}
    h^{\mathrm{expr}}_i &= f_{\mathrm{expr}}(\log(1+x_i)), \\
    h^{\mathrm{mask}}_i &= f_{\mathrm{mask}}(m_i), \\
    h_i &= \gamma(h^{\mathrm{mask}}_i) \odot h^{\mathrm{expr}}_i
           + \beta(h^{\mathrm{mask}}_i),
\end{align}
where \(h_i\) is a deterministic hidden representation, and
\(\gamma(\cdot)\) and \(\beta(\cdot)\) are learned modulation functions. This design
allows the encoder to distinguish true zeros from unmeasured genes and to adapt the
representation to the available panel.

\subsection{Gaussian-mixture latent prior}

To represent multiple cellular states in a generative latent space, we use a Gaussian
mixture prior over the latent variable \(z_i \in \mathbb{R}^{D}\):
\begin{equation}
    p(z_i) = \sum_{k=1}^{K} \pi_k
    \mathcal{N}\left(z_i \mid \mu_k, \Sigma_k\right),
\end{equation}
where \(K\) is the number of mixture components, \(\pi_k\) are learned mixture
weights, \(\mu_k \in \mathbb{R}^{D}\) are component centers, and \(\Sigma_k\) are
component-specific covariance matrices. In the main experiments described here,
\(\Sigma_k\) is diagonal unless otherwise stated. Thus, for \(K=64\) and
\(D=256\), the prior contains 64 prototype centers in a 256-dimensional latent
space. The mixture components are shared across tissues by default; optional
conditioning can be enabled for selected decoder covariates, as described below.

\subsection{Mixture posterior}

The posterior is parameterized as a Gaussian-mixture variational posterior with a
discrete component variable \(c_i\):
\begin{equation}
    q_{\phi}(z_i,c_i \mid x_i,m_i)
    = q_{\phi}(c_i \mid h_i)\,
      q_{\phi}(z_i \mid h_i,c_i).
\end{equation}
The encoder hidden state \(h_i\) is used to predict a categorical distribution over
components,
\begin{equation}
    q_{\phi}(c_i=k \mid h_i)
    = \mathrm{softmax}_k(a_{\phi}(h_i)),
\end{equation}
and component-specific posterior parameters:
\begin{equation}
    q_{\phi}(z_i \mid h_i,c_i=k)
    =
    \mathcal{N}\left(
        z_i \mid \mu_{\phi k}(h_i), \mathrm{diag}(\sigma^2_{\phi k}(h_i))
    \right).
\end{equation}
During training, the component assignment is relaxed with a differentiable soft
selection, and during evaluation the most likely component can be used for sampling.
For deterministic embedding extraction, we use either the encoder hidden state
\(h_i\), the base posterior mean, or the mixture-averaged posterior mean
\(\sum_k q_{\phi}(c_i=k \mid h_i)\mu_{\phi k}(h_i)\), depending on the analysis.

\subsection{Count decoder}

The decoder maps \(z_i\) back to a full-gene expression distribution. It predicts gene
logits and a cell-level library size:
\begin{align}
    r_i &= f_{\theta}(z_i, b_i), \\
    \ell_i &= \mathrm{softplus}(g_{\theta}(z_i)), \\
    \rho_{ig} &= \mathrm{softmax}(r_i)_g, \\
    \lambda_{ig} &= \ell_i \rho_{ig}.
\end{align}
Here \(b_i\) denotes an optional batch condition. In the experiments using batch
conditioning, \(b_i\) is provided to the decoder through FiLM-style modulation, while
tissue conditioning is disabled unless explicitly stated. This allows the decoder to
explain donor- or batch-specific expression shifts while preserving an unconditioned
latent representation. To prevent the decoder from relying exclusively on batch
metadata, an additional batchless reconstruction loss can be computed with the batch
condition set to zero.

For count likelihoods, we use either a Poisson model,
\begin{equation}
    x_{ig} \sim \mathrm{Poisson}(\lambda_{ig}),
\end{equation}
or a negative binomial model,
\begin{equation}
    x_{ig} \sim \mathrm{NB}(\lambda_{ig}, \theta_g),
\end{equation}
where \(\theta_g\) is a learned gene-level dispersion parameter. Unless otherwise
specified, the likelihood is evaluated only for observed genes:
\begin{equation}
    \mathcal{L}_{\mathrm{rec}}
    =
    - \frac{1}{\sum_g m_{ig}}
      \sum_{g=1}^{G} m_{ig}
      \log p_{\theta}(x_{ig} \mid z_i).
\end{equation}

\subsection{Variational objective}

The core variational objective combines observed-gene reconstruction with the
Gaussian-mixture KL divergence:
\begin{equation}
    \mathcal{L}_{\mathrm{VAE}}
    =
    \mathcal{L}_{\mathrm{rec}}
    +
    \beta\,
    \mathrm{KL}\left[
        q_{\phi}(z_i,c_i \mid x_i,m_i)
        \,\|\, p_{\theta}(z_i,c_i)
    \right].
\end{equation}
The KL term decomposes into a continuous component and a discrete component:
\begin{align}
    \mathrm{KL}
    &=
    \sum_{k=1}^{K} q_{\phi}(c_i=k \mid h_i)
    \mathrm{KL}\left[
        q_{\phi}(z_i \mid h_i,c_i=k)
        \,\|\, p_{\theta}(z_i \mid c_i=k)
    \right] \nonumber \\
    &\quad+
    \mathrm{KL}\left[
        q_{\phi}(c_i \mid h_i)
        \,\|\, p_{\theta}(c_i)
    \right].
\end{align}
The KL weight \(\beta\) is kept small in our main experiments to preserve biological
variation in the embedding while maintaining a usable generative prior.

\subsection{Panel-view contrastive learning}

Because the same biological cell may be observed through different gene panels, we
introduce a panel-view contrastive objective. For each training cell, we construct two
views: a real view using the original observed mask and a masked view generated by
randomly hiding a subset of already observed genes. The two views share the same count
vector but differ in encoder-visible masks. Let \(u_i\) and \(v_i\) denote the
deterministic embeddings of the real and masked views, respectively. In the main
configuration, these embeddings are the encoder hidden states \(h_i\). We use a
symmetric InfoNCE loss:
\begin{equation}
    \mathcal{L}_{\mathrm{con}}
    =
    \frac{1}{2}
    \left[
    \mathrm{CE}
    \left(
        \frac{UV^{\top}}{\tau}, \mathrm{diag}
    \right)
    +
    \mathrm{CE}
    \left(
        \frac{VU^{\top}}{\tau}, \mathrm{diag}
    \right)
    \right],
\end{equation}
where rows of \(U\) and \(V\) are normalized embeddings from the two views, and
\(\tau\) is the contrastive temperature. This objective directly encodes the desired
panel invariance: two partial observations of the same cell should map to nearby
representations.

\subsection{Auxiliary biological supervision}

When cell-type annotations are available, we add an auxiliary classification head on
the latent representation:
\begin{equation}
    \mathcal{L}_{\mathrm{ct}}
    =
    \mathrm{CE}\left(s_{\theta}(z_i), y_i\right),
\end{equation}
where labels with missing annotation are ignored. This term is used as a weak semantic
anchor rather than as the primary training objective.

\subsection{Reweighting rare gene and cell signals}

To reduce domination by highly expressed genes, reconstruction losses can be
gene-weighted. We maintain exponential moving averages of observed-gene means and
second moments and compute a coefficient-of-variation statistic,
\begin{equation}
    \mathrm{CV}_g =
    \frac{\sqrt{\mathbb{E}[x_g^2] - \mathbb{E}[x_g]^2}}
         {\mathbb{E}[x_g] + \epsilon}.
\end{equation}
Gene weights are then formed as
\begin{equation}
    w_g = 1 + \log(1+\mathrm{CV}_g),
\end{equation}
normalized, clipped, and mixed with a uniform weight. The final reconstruction term is
multiplied by \(w_g\) at observed gene entries. We also optionally upweight cells in
under-represented local structures by running a lightweight within-batch clustering
procedure and assigning larger reconstruction weights to cells in smaller clusters.

\subsection{Full training objective}

The full objective is
\begin{equation}
    \mathcal{L}
    =
    \mathcal{L}_{\mathrm{VAE}}
    +
    \lambda_{\mathrm{con}}\mathcal{L}_{\mathrm{con}}
    +
    \lambda_{\mathrm{ct}}\mathcal{L}_{\mathrm{ct}}
    +
    \lambda_{\mathrm{batchless}}\mathcal{L}_{\mathrm{batchless}}.
\end{equation}
The relative weights control the balance between generative modeling, panel-invariant
representation learning and weak biological supervision.
